\subsection{量词的性质}

从现在起，我们将只考虑带量词的理论。在本节的其余部分， \(\mathfrak{C}\) 将表示这样的一个理论，而 \(\mathfrak{C}_0\) 将表示没有显式公理、具有与 \(\mathfrak{C}\) 相同符号且其唯一方案为S1至S5的理论。 \(\mathfrak{C}\) 强于 \(\mathfrak{C}_0\) .

C29. 设R为C中的一个关系，并设x为一个字母。则关系“非 \((\exists x)R\) ”和 \((\forall x)\) (非 R) 在 C 中等价。

只需在理论 \(\mathfrak{G}_0\) 中证明该判据即可，在其中 \(x\) 不是常量。定理 \(R \Longleftrightarrow\) (非非 \(R\) ) 通过 C3（ \(\S 2\) ，第 3 号）给我们定理

\[
(\exists x) R \Longrightarrow (\tau_ {x} (R) | x) (\text { 非   非 } R)
\]

并且

\[
(\exists x) (\text { 非   非 } R) \Longrightarrow (\tau_ {x} (\text { 非   非 } R) | x) R.
\]

应用 S5，我们推导出在 \(\mathfrak{C}_{0}\)

\[
(\exists x) R \Longrightarrow (\exists x) (\text { 非   非 } R), \quad (\exists x) (\text { 非   非 } R) \Longrightarrow (\exists x) R,
\]

中的定理，由此得到定理 \((\exists x)R \Longleftrightarrow (\exists x)\) (非非 \(R\) ）。现在

\[
(\exists x) (\text { 非   非 } R)
\]

在 \(\mathcal{G}_0\) 中等价于“非非 \((\exists x)\) (非非 \(R\) ）”，即“非 \((\forall x)\) （非 \(R\) ）”。因此结论成立。

准则C28和C29使我们能够从其中一个量词的性质推导出另一个量词的性质。

C30. 设R是G中的一个关系，T是G中的一个项，且x是一个字母。则关系 \((\forall x)R \Longrightarrow (T|x)R\) 是G中的一个定理。

根据S5， \((T|x)\) （非 \(R\) ) \(\Longrightarrow\) ( \(\tau_{x}\) （非 \(R\) ) \(|x\) )（非 \(R\) ）是一条公理。这个关系与下式相同：

\[
(\text { 非 } (T | x) R) \Longrightarrow \text { 非 } ((\tau_ {x} (\text { 非 } R) | x) R).
\]

因此， \((\tau_{\mathbf{x}}(\text{非 } R)|x)R \Longrightarrow (T|x)R\) 是 \(\mathfrak{C}\) 中的定理 。现在使用C26（第1号）。

设 \(\pmb{R}\) 是 \(\mathfrak{C}\) 中的一个关系。根据C26、C27和C30（前提是字母 \(\pmb{x}\) 不是 \(\mathfrak{C}\) 的常量），无论我们陈述定理 \(\pmb{R}\) 于 \(\mathfrak{C}\) 中，还是陈述定理 \((\forall x)\pmb{R}\) ，还是陈述元数学规则：如果 \(\pmb{T}\) 是 \(\mathfrak{C}\) 中的任意项，那么 \((\pmb{T}|\pmb{x})\pmb{R}\) 是 \(\mathfrak{C}\) 中的定理 ，结果都是相同的。

C31. 设R和S是C中的关系，并设x是一个不是C的常量的字母。如果 \(R \Longrightarrow S\) （相应地， \(R \Longleftrightarrow S\) ）是 C 中的一个定理，那么

\[
\begin{array}{c c} (\forall x) R \Longrightarrow (\forall x) S, & (\exists x) R \Longrightarrow (\exists x) S \\ \text{[分别 } (\forall x) R \Longleftrightarrow (\forall x) S, & (\exists x) R \Longleftrightarrow (\exists x) S \text{]} \end{array}
\]

是 \(\mathfrak{G}\) 中的定理 .

假设 \(R \Longrightarrow S\) 是 \(\mathcal{G}\) 中的定理 。让我们附加假设 \((\forall x)R\) （其中 \(x\) 不出现）。那么 \(R\) ，因此 \(S\) ，因此也有 \((\forall x)S\) ，是真的。因此 \((\forall x)R \Longrightarrow (\forall x)S\) 是 \(\mathcal{G}\) 中的定理 。由此可知，如果 \(R \Longleftrightarrow S\) 是 \(\mathcal{G}\) 中的定理 ，那么也有

\[
(\forall x) R \Longleftrightarrow (\forall x) S.
\]

关于 \(\exists\) 的规则现在可以利用C29推导出来。

C32. 设R和S是C中的关系，并设x是一个字母。则关系

\[
\begin{array}{l} (\forall x) (R \text {   且   } S) \iff ((\forall x) R \text {   且   } (\forall x) S), \\ (\exists x) (R \text {   或   } S) \iff ((\exists x) R \text {   或   } (\exists x) S) \end{array}
\]

是 \(\mathfrak{C}\) 中的定理 .

只需在 \(C_{0}\) 中证明这些准则即可，其中x不是常量。如果 \((\forall x)(R \text{ 且 } S)\) 成立，则``R 且 \(S\)''为真，因此关系R、S各自为真。从而每个关系 \((\forall x)R\) , \((\forall x)S\) 是真的，因此`` \((\forall x)R\) 和 \((\forall x)S\)''为真。类似地，可以证明如果`` \((\forall x)R\) 和 \((\forall x)S\)'' 为真，则 \((\forall x)(R \text{ 且 } S)\) 成立。因此第一个定理成立。第二个定理通过应用C29得出。

需要注意的是，如果 \((\forall x)(R \text{ 或 } S)\) 是G中的一个定理，我们不能得出结论说 \(((\forall x)R \text{ 或 } (\forall x)S)\) 是G中的一个定理。直观地说，说这个关系 \((\forall x)(R \text{ 或 } S)\) 为真，意味着对于每个对象x，关系R和S中至少有一个成立；但一般而言，两者中只有一个会成立，至于究竟是R还是S，则取决于x的选择。同样地，如果 \(((\exists x)R \text{ 且 } (\exists x)S)\) 是G中的一个定理，我们不能得出结论说 \((\exists x)(R \text{ 且 } S)\) 是 G 中的一个定理。然而，有以下判据：

C33. 设 R 和 S 是 G 中的关系，并设 x 是不出现在 R 中的一个字母。那么关系

\[
\begin{array}{l l} (\forall x) (R \text {   或   } S) & \Longleftrightarrow (R \text {   或   } (\forall x) S), \\ (\exists x) (R \text {   且   } S) & \Longleftrightarrow (R \text {   且   } (\exists x) S) \end{array}
\]

是 \(\mathfrak{C}\) 中的定理 .

在 \(\mathfrak{C}_0\) 中建立该判据就足够了，其中 \(x\) 不是常量。设 \(\mathfrak{C}'\) 为通过将 \((\forall x)(R\) 或 \(S)\) 附加到 \(\mathfrak{C}_0\) 的公理而得到的理论。在 \(\mathfrak{C}'\) 中，`` \(R\) 或 \(S\) ''以及因此（非 \(R) \Longrightarrow S\) ，都是定理。如果``非 \(R\) ''为真（一个不涉及 \(x\) 的假设），那么 \(S\) 以及因此 \((\forall x)S\) 都为真。从而

\[
(\text { 非 } R) \Longrightarrow (\forall x) S
\]

是 \(\mathcal{C}'\) 中的定理 ，因此 \((\forall x)(R\) 或 \(S) \Longrightarrow (R\) 或 \((\forall x)S)\) 是 \(\mathcal{C}_0\) 中的定理 。同样地，如果`` \(R\) 或 \((\forall x)S\) '' 为真，则 `` \(R\) 或 \(S\) '' 以及因此 \((\forall x)(R\) 或 \(S)\) 都为真。从而

\[
(R \text {   或   } (\forall x) S) \Longrightarrow (\forall x) (R \text {   或   } S)
\]

是 \(\mathfrak{G}_0\) 中的定理。关于 \(\exists\) 的规则可通过应用C29得出。

C34. 设R为一个关系，并设x和y为字母。则关系

\[
\begin{array}{l} (\forall x) (\forall y) R \iff (\forall y) (\forall x) R, \\ (\exists x) (\exists y) R \iff (\exists y) (\exists x) R, \\ (\exists x) (\forall y) R \implies (\forall y) (\exists x) R \end{array}
\]

是 \(\mathfrak{G}\) 中的定理 .

只需在 \(\mathcal{G}_0\) 中建立该判据就足够了，其中 \(x\) 和 \(y\) 不是常量。如果 \((\forall x)(\forall y)R\) 成立，则 \((\forall y)R\) ，从而 \(R\) ，因此 \((\forall x)R\) ，因此 \((\forall y)(\forall x)R\) ，都成立。同样地，如果 \((\forall y)(\forall x)R\) 成立，则 \((\forall x)(\forall y)R\) 是真的；第一个定理由此得证。第二个定理现在可通过使用C29得出。最后，由于 \((\forall y)R \Longrightarrow R\) 是 \(\mathcal{C}_0\) 中的定理 ，所以 \((\exists x)(\forall y)R \Longrightarrow (\exists x)R\) 由C31可知；如果 \((\exists x)(\forall y)R\) 成立，则 \((\exists x)R\) 成立，因此 \((\forall y)(\exists x)R\) 也成立。因此第三个定理成立。

另一方面，如果 \((\forall y)(\exists x)R\) 是 C 中的定理，我们不能推出 \((\exists x)(\forall y)R\) 是 C 中的定理。直观地说，关系 \((\forall y)(\exists x)R\) 为真意味着：给定任意对象 y，存在对象 x 使得 R 是对象 x 与 y 之间的一个真关系。但一般而言，对象 x 依赖于对象 y 的选择，而说 \((\exists x)(\forall y)R\) 为真则意味着存在一个固定的对象 x，使得 R 是该对象与任意对象 y 之间的真关系。

